\documentclass[11pt]{article}
\usepackage[margin=1in]{geometry}
\usepackage{amsmath,amssymb,amsthm}
\usepackage[T1]{fontenc}
\usepackage{lmodern}
\usepackage{microtype}
\usepackage{booktabs}
\usepackage[title]{appendix}
\usepackage{xurl}
\usepackage{hyperref}
\hypersetup{colorlinks=true,linkcolor=blue,citecolor=blue,urlcolor=blue,
  pdftitle={A placebo-controlled test of WSPR signal anomalies on radio paths crossed by tracked aircraft},pdfauthor={Peter Graham}}
\title{A placebo-controlled test of WSPR signal anomalies on radio paths crossed by tracked aircraft}
\author{Peter Graham\\ \small\url{https://orcid.org/0009-0007-8493-1505}}
\date{November 23, 2026}
\pagestyle{plain}
\DeclareUnicodeCharacter{0394}{\ensuremath{\Delta}}
\begin{document}
\maketitle
\begin{abstract}
Godfrey and colleagues proposed that aircraft near an amateur-radio signal path make unusual WSPR reception reports more frequent. They used such reports to propose positions for flight MH370. We tested an adapted anomaly flag using aircraft positions broadcast by ADS-B. We compared reports on covered paths (recorded traffic in every grid cell) passing within 5~km of a tracked aircraft at or above 20{,}000~ft with matched reports on covered paths away from tracked aircraft. We repeated the comparison using the previous and next sampled dates' aircraft. The observed difference in flag rates minus the mean of these placebo differences is Δ. Across two samples of 36 dates, Δ was $+0.21$ percentage points with same-link matching (95\% interval $-0.53$ to $0.95$) and $+0.35$ with shared-station matching ($-0.46$ to $1.17$). The date-bootstrap intervals are nominal because dates share aircraft data. Both upper limits stayed below the prespecified two-point threshold after any single-date deletion. Both designs detected a simulated 6~dB signal-to-noise increase on a random 10\% of all records crossed under the real positions before coverage filtering and matching. Only 0.3\% of covered crossed records had a match. We selected the follow-up after the first result and before downloading its data.

This work has not yet been reviewed by experts.
The analysis, the code and the text were produced by AI models from OpenAI and Anthropic, working under the author's direction.
\end{abstract}
\medskip
\noindent\textbf{Keywords:} WSPR, amateur radio, ADS-B, MH370, placebo control, pre-specified analysis.

\section{Introduction}

\subsection{Main result}

Godfrey and colleagues proposed that an aircraft crossing a WSPR signal path raises the rate of anomalous reports. They used anomalies to propose positions for flight MH370~\cite{oefgr,mh370fpa}. Does an adapted flag occur more often near tracked aircraft than on matched paths away from them, beyond what we would see with other days' aircraft? An increase would support this prediction. Our decision rule requires both upper limits to fall below two percentage points and both comparisons to detect a simulated increase before we call the outcome negative.

WSPR (Weak Signal Propagation Reporter) is an amateur-radio protocol for weak transmissions lasting just under two minutes~\cite{taylor}. Receiving stations configured to upload send their decoded reports to WSPRnet. Each gives the time, station identifiers and grid locations, band, received frequency, signal-to-noise ratio (SNR) and frequency drift. ADS-B Exchange archives aircraft position broadcasts~\cite{adsbx}. A \emph{link} pairs a transmitter and receiver on one band. A \emph{link-slot} is a reception record for that link in a two-minute slot. We approximate the path by the great-circle arc between the centers of the stations' grid squares (six-character Maidenhead locators, accurate to about 5~km).

On covered paths, where every sampled $1^\circ\times1^\circ$ grid cell has regular aircraft reports, 96.4\% of 26{,}692{,}421 scored records were crossed. The flag marked 67.6\% of crossed records and 63.7\% of eligible uncrossed records. With aircraft and flags both common, a coincidence is easy to find~\cite{loi}. Only 96{,}589 records qualified as uncrossed controls, leaving matches for about 0.3\% of the 25{,}741{,}283 crossed records. Counting each record once, median path lengths were 958~km for covered crossed records and 197~km for eligible uncrossed records. For matched same-link cases on the first 36 dates, the median of the daily median lengths was 271~km.

For each crossed case we subtract the flagged share of its controls from its own flag, then average over cases. We repeat this comparison with the next and previous sampled dates' aircraft positions. The real average minus the mean of those two placebo averages is $\Delta$, reported in percentage points (extra flags per 100 matched cases).

\begin{quote}
\textbf{Result.} For the 72 sampled dates, $\Delta$ was $+0.21$ points with same-link matching (74{,}736 cases, nominal 95\% interval $-0.53$ to $0.95$) and $+0.35$ with shared-station matching (67{,}253 cases, $-0.46$ to $1.17$). The simulated SNR increase gave $+2.96$ ($2.19$ to $3.70$) and $+2.68$ ($1.77$ to $3.60$) points, respectively. The pooled data meet the prespecified rule for a negative outcome.
\end{quote}

\subsection{Prior work}\label{sec:prior}

Westphal proposed using archived WSPR data to investigate MH370 in 2021 and compared signals between Antarctica and New Zealand with Flightradar24 positions~\cite{westphal2021}. Godfrey and colleagues tested individual flights against absence comparisons. For OE-FGR, Godfrey and Coetzee compared the same links two, four or six minutes away. The area under the receiver operating characteristic (ROC) curve was 59\% over 183 observations (50\% is chance discrimination)~\cite{oefgrpdf}. For QTR901, Godfrey, Coetzee and Maskell compared links within one nautical mile (our sensitivity radius) with nearby links and with observations six minutes later. Their area was 58\% for ``244 observations during all phases of the flight'' (Fig.~22). Page~17 places some phases outside ADS-B range~\cite{qtr901pdf}. Their 48-flight study reported 55.6\% over 1{,}078 observations using a 2~km radius and an absence comparison six minutes later with no aircraft within that radius~\cite{gpr}. Other reports cover short paths, THY161 and a proposed MH370 track~\cite{wsprradar,shortdist,thy161,mh370fpa}.

Anderson and the Liverpool team both considered all traffic on a day. In December 2021 Anderson binned WSPR frequency offsets by distance to aircraft at or above 20{,}000~ft and called for more than one day of data~\cite{anderson}. We use Anderson's altitude threshold and extend the comparison to 72 dates. The University of Liverpool announced a study of WSPR and known flight paths~\cite{liverpool}. Godfrey described Maskell's ROC comparison of all airborne aircraft in the ADS-B data against all WSPR data on May 1, 2023, one of our first-sample dates~\cite{liverpooldesign}. We found no published outcome of that comparison.

Loi and colleagues checked WSPR detections against Flightradar24 positions within 5~km of the transmitter--receiver line, the radius we use with ADS-B. They wrote: ``We are unable to conclusively verify that data from WSPRnet can be used to detect and localize aircraft. Our preliminary verification suggests that using standard score to find anomalies in SNR could be used to detect a aircraft flying through a WSPR link.'' They cautioned that many intersections and aircraft allow chance matches~\cite{loi}.

Marten's WSPR Tracker compared anomaly rates on the same links at ``the same time of day, one day earlier (no flight present)''~\cite{martencontrol}. Marten's localization backtest on two held-out flights gave a mean of 4{,}631.8~km for the flights' median position errors~\cite{marten}. We measure a difference in flag rates with matched uncrossed controls, holding the radio records fixed and shifting the aircraft positions.

Kent argued that a scattered signal would be too weak to detect~\cite{kent}. Iannello recorded protocol author Joe Taylor's objection to tracking aircraft with historical WSPR data. Gilbert reported that the group checked the QF6036 blind test's nine position indicators and found none correct among the seven with ADS-B data~\cite{iannello2021}.

\section{Data and method}\label{sec:method}

The release gives the protocol and its SHA-256 at the time the analysis was specified. The author's commit history records that the protocol preceded the first outcome calculation and the follow-up sampling rule preceded its data download. There is no independent timestamp for that order. Appendices~\ref{app:defs} and~\ref{app:dates} give definitions and dates.

\paragraph{Data.}\label{sec:data} WSPR records come from the wspr.live mirror of WSPRnet~\cite{wsprlive}, from 21:00 UTC on the day before each sampled date to 01:00 UTC on the day after. ADS-B Exchange makes complete first-of-month data free and has archived snapshots every five seconds since May 2020. We select one snapshot a minute, 488 per date~\cite{adsbx}. We use the first hour of each of the eight three-hour periods starting at 00, 03, \ldots, 21 UTC, giving 30 two-minute test slots per period. The first sample is 36 first-of-month dates from December 2022 to December 2025. The second, fixed after the first sample's outcome and before its data were downloaded, is the 36 most recent unused first-of-month dates with complete snapshots, from August 2020 to November 2026.

\paragraph{The anomaly rule.}\label{sec:rule} We test an adaptation of the rule Godfrey stated in 2024 comments on the 2022 OE-FGR case study~\cite{oefgr}; Appendix~\ref{app:rule} maps each phrase. A record is flagged when its SNR or frequency lies more than 0.75 sample standard deviations from the mean of the same link's records within $\pm3$~hours of it, or when its drift is nonzero and fewer than half of those records drift. Records are grouped by transmitter, receiver and band, where Godfrey's triple uses frequency. A record with fewer than five records in its window gets no flag at all and is not scored; Godfrey instead extends sparse windows up to $\pm3$~days. Stations with a slash in the call sign, balloon-telemetry call signs, or more than one locator string among all downloaded records are dropped as possibly mobile, before only six-character locators are kept. Paths shorter than 100~km are dropped.

\paragraph{Crossings.} For each two-minute slot we reconstruct tracks from three one-minute snapshots. A track counts as high for the slot if any snapshot puts it at or above 20{,}000~ft. A path is crossed if a high track comes within 5~km of its arc. A control must have every tracked high aircraft at least 25~km away and every tracked airborne aircraft at least 5~km away (Appendix~\ref{app:defs}). The \emph{coverage mask} keeps a path when every $1^\circ\times1^\circ$ cell along its arc holds aircraft reports, airborne or on the ground, in at least a quarter of that period's snapshots. The mask is a traffic proxy, not proof that every aircraft present was recorded, so ``uncrossed'' means that no recorded track came near.

\paragraph{Designs, placebos and uncertainty.} \emph{Same-link matching} takes as controls the same link's eligible uncrossed link-slots on the same date with report times within 30 minutes of the case. \emph{Shared-station matching} takes eligible uncrossed link-slots with the same report time and band, on another link sharing the case's transmitter or its receiver in the same role, with path length 0.8 to 1.25 times the case's. Matching in time and by station is meant to reduce, not remove, differences in propagation conditions. For each case, $d$ is its flag minus the share of its controls that are flagged, and $E$ is the mean of $d$ over the cases; a flagged case with four controls, one of them flagged, gives $d=1-\tfrac14=0.75$. We then replace the aircraft positions by those of the next and of the previous date of the same sample at the same UTC times, wrapping from the last date to the first, keep the radio records as they are, and redo the labels, the matching and $E$; the shifted positions label different records as crossed, so each average uses its own cases. In proportion units the primary statistic is
\[
  \Delta \;=\; E_{\text{real}}-\tfrac12\bigl(E_{\text{next}}+E_{\text{prev}}\bigr).
\]
The shifted positions keep the daily pattern of traffic but break the timing between a record and the recorded aircraft, so the shifted $E$ values show how large an excess the selection and matching produce without that timing. Section~\ref{sec:limits} states what an aircraft-effect reading would need. Intervals are percentiles of 10{,}000 resamples of whole sampled dates, each date's real and shifted sums resampled together. Dates share aircraft data through the shifted positions, so the intervals are nominal.

\paragraph{Simulated increase.} In a copy of the records, each link-slot crossed under the real positions gets $+6$~dB of SNR with probability 10\%, before coverage filtering and matching (pooled, 4{,}128{,}269 of 41{,}301{,}615; 2{,}572{,}854 of them crossed and covered). Every record's SNR flag is recomputed on the copy, so neighboring records' windows change too, and the real and shifted analyses are rerun on it.

\paragraph{Decision rule.} In order: if either simulated-increase lower limit is not above zero, the outcome is inconclusive; if both primary lower limits are above zero, it is positive; if both primary upper limits are below 2 percentage points, it is negative; otherwise it is inconclusive. Against flag rates near 65\%, the protocol treated a difference under two percentage points as too small to matter for this test. It is a decision threshold, not an established limit on any tracking method.

\section{Results}\label{sec:results}

\begin{table}[h]
\centering
\small
\setlength{\tabcolsep}{3pt}
\begin{tabular}{llrrrl}
\toprule
Sample & Matching & Cases & $E_{\text{real}}$ & $E_{\text{next}}$, $E_{\text{prev}}$ & $\Delta$ (95\% interval) \\
\midrule
First (36) & Same-link & 36{,}359 & $+0.47$ & $-0.19$, $-0.37$ & $+0.75$ ($-0.37$, $1.85$) \\
 & Shared-station & 32{,}371 & $+0.69$ & $+0.06$, $+0.89$ & $+0.22$ ($-0.95$, $1.29$) \\
Second (36) & Same-link & 38{,}377 & $-0.24$ & $-0.02$, $+0.14$ & $-0.29$ ($-1.22$, $0.65$) \\
 & Shared-station & 34{,}882 & $+0.31$ & $+0.08$, $-0.38$ & $+0.46$ ($-0.64$, $1.65$) \\
Pooled (72) & Same-link & 74{,}736 & $+0.11$ & $-0.10$, $-0.10$ & $+0.21$ ($-0.53$, $0.95$) \\
 & Shared-station & 67{,}253 & $+0.49$ & $+0.07$, $+0.21$ & $+0.35$ ($-0.46$, $1.17$) \\
\midrule
Simulated & Same-link & & & & $+2.96$ ($2.19$, $3.70$) \\
 & Shared-station & & & & $+2.68$ ($1.77$, $3.60$) \\
\bottomrule
\end{tabular}
\caption{$E$ and $\Delta$ in percentage points, for each sample and pooled. Cases are crossed link-slots with at least one control under the real positions. The simulated-increase rows are $\Delta$ after adding 6~dB of SNR with probability 10\% to every link-slot crossed under the real positions, before coverage filtering and matching (Section~\ref{sec:method}); for each sample alone they are in Appendix~\ref{app:variants}.}\label{tab:main}
\end{table}

Table~\ref{tab:main} gives the real and placebo comparisons separately for each sample and for the pooled data.

\paragraph{Why a second sample was added.} On the first 36 dates alone, the same-link upper limit was 1.85. Removing December 1, 2025 raised it to 2.01, past the threshold; on that date the same-link real excess is $-2.61$ points over 1{,}441 cases, against $+0.39$ and $+0.82$ with the shifted positions. Because one date decided the outcome, the second sample (Section~\ref{sec:data}) and a robustness condition were fixed before its data were downloaded: the negative outcome counts as robust only if the pooled intervals meet the rule under every single-date deletion and every block length below.

\paragraph{Date deletions.} Every single-date deletion kept both pooled upper limits below two points. Dropping a date's sums gave maximum upper limits of 1.02 points (same-link) and 1.25 (shared-station). Also withholding that date's aircraft from the two neighboring placebo comparisons gave maxima of 1.06 and 1.35. The lowest simulated-increase lower limit across both deletion methods was 1.60.

\paragraph{Blocks of dates.} Resampling blocks of 2, 3, 4 or 6 successive dates also preserved the pooled outcome. Across block lengths the outermost interval endpoints were $-0.58$ and $0.99$ points for same-link matching and $-0.48$ and $1.15$ for shared-station matching. Chronological blocks preserve some dependence from reused aircraft, but shifted positions wrap within each sample.

\paragraph{Matching choices.} Matching on slot starts, dropping odd-minute reports at matching or making the control cutoffs strict, separately or together, moved $\Delta$ by at most 0.02 points (Appendix~\ref{app:variants}).

\paragraph{Variants fixed in advance.} None decides the outcome (Appendix~\ref{app:variants}). Pooled, the 22 upper limits of the binary-flag variants are below 2 percentage points (the largest $1.56$), and none of the 24 variant intervals excludes zero. In each sample alone a different one excludes zero: the same-link drift flag in the first sample ($+0.43$, $0.07$ to $0.80$) and the shared-station SNR flag in the second ($+1.21$, $0.05$ to $2.36$). We treat both as exploratory.

\paragraph{Crossing geometry.} On 960 sampled first-sample link-slots near the cutoffs or used as controls, a reference calculation using every aircraft in the raw snapshots changes three labels, one of them on an odd-minute record, and otherwise agrees within 0.5~km (Appendix~\ref{app:checks}).

\section{Scope and interpretation}\label{sec:limits}

\paragraph{What is bounded.} The bound applies to the matched cases. The available same-link length profile favors paths of a few hundred kilometers. The pooled average does not bound an effect on long paths separately, on crossings without controls, or near unrecorded aircraft. Interpreting $\Delta$ as an aircraft effect would require the real and shifted positions to select comparable paths and the controls to share the cases' propagation conditions.

\paragraph{Coverage and geometry.} The mask keeps only paths whose cells have recorded traffic, so the result says nothing about paths it removes, many of them over oceans. The geometry is that of locator centers and surface arcs; the radio signal travels by the ionosphere and can follow a different route.

\paragraph{Uncertainty.} Block and deletion checks probe the dependence between dates that share aircraft data. The pooled interval makes no allowance for the decision to extend the study after seeing the first result. The second sample, specified before downloading its data, meets the rule on its own. The simulation checks sensitivity to one particular increase.

\paragraph{What the result suggests.} One possibility is that this flag carries little information about tracked aircraft near a path. Another, as Loi and colleagues cautioned~\cite{loi}, is that the large numbers of aircraft and anomalous paths make apparent matches common by chance.

\section{Verification and reproducibility}\label{sec:verify}

Two independent implementations were written from the protocol without our code. The first, on 36 dates, reproduced the outcome (same-link $+0.54$ points, shared-station $-0.09$) but matched 21--26\% more cases and controls, outside the 5\% tolerance. On February 1, 2023 its locator filter kept 26.5\% more records than ours, close to that date's 26\% gap in scored records. Other differences were not harmonized. The second implementation used the protocol and addendum on all 72 dates. It reproduced the pooled outcome (same-link $+0.21$ points, shared-station $+0.32$), passed every deletion and block check, and matched case and control counts within 5\%.

Exact arithmetic confirmed the flags at the 0.75-standard-deviation threshold for all 45{,}033{,}331 scored records, including simulated SNR. The 5{,}633 SNR and 20{,}631 frequency scores exactly at the threshold were unflagged, as specified.

The protocol, addendum, code, results and replay instructions accompany version~3 at
\begin{center}\small\url{https://github.com/6/clank-lab/tree/wspr-aircraft-anomalies-v3/results/wspr-aircraft-anomalies}\end{center}
The scripts download the source data.

\section*{Acknowledgments}
The hypothesis and the anomaly rule are Richard Godfrey's. The OE-FGR report is by Godfrey and Hannes Coetzee, and the QTR901, THY161, short-distance, MH370 and global passive radar reports by Godfrey, Coetzee and Simon Maskell. Robert Westphal proposed using archived WSPR data for MH370 in 2021~\cite{westphal2021}. The University of Liverpool team (Simon Maskell, Alberto Acuto, Paul Horridge and Shashi Lakra) announced their larger study. The critiques of Victor Iannello, Steve Kent, Mick Gilbert and others, and Ed Anderson's one-day comparison, shaped the controls. The data are from the WSPRnet community through wspr.live and from ADS-B Exchange.

\begin{thebibliography}{99}
\bibitem{adsbx} ADS-B Exchange, sample data, \url{https://www.adsbexchange.com/data-products/sample-data/}, and \texttt{readsb-hist} snapshots, \url{https://samples.adsbexchange.com/}.
\bibitem{anderson} E.~Anderson (370Location), comment of December 28, 2021 on~\cite{iannello2021}, \url{https://mh370.radiantphysics.com/2021/12/19/wspr-cant-find-mh370/#comment-32855}.
\bibitem{gpr} R.~Godfrey, H.~Coetzee and S.~Maskell, A proposed global passive radar, January 1, 2025, \url{https://www.dropbox.com/scl/fi/8q54o5btw9mfkc6io30ye/Global-Passive-Radar-01JAN2025.pdf?rlkey=kk4xbt04c8tlojs3j294lpljw&dl=0}, linked from \url{https://www.mh370search.com/2025/01/01/new-technology/}; pp.~2, 6 and 16.
\bibitem{mh370fpa} R.~Godfrey, H.~Coetzee and S.~Maskell, MH370 flight path analysis: case study, August 31, 2023, \url{https://www.dropbox.com/scl/fi/nn3eedtd9ew15f81n2xx7/MH370-Flight-Path-Analysis-31st-August-2023.pdf?rlkey=c72bo3kc2vv74swr3wll3e4j9&dl=0}, listed at \url{https://www.mh370search.com/archive/}; \S\S6 and 9.
\bibitem{shortdist} R.~Godfrey, H.~Coetzee and S.~Maskell, How does WSPR detect aircraft over short distances? Technical paper, December 17, 2023, \url{https://www.dropbox.com/scl/fi/8uz3l31ayfpbjjo7q3vwu/How-Does-WSPR-Detect-Aircraft-over-Short-Distances-17DEC2023.pdf?rlkey=bk2enssvhdv58xa0viuzpupv0&dl=0}; pp.~1 and 61.
\bibitem{thy161} R.~Godfrey, H.~Coetzee and S.~Maskell, How does WSPR detect aircraft over long distances? Technical paper, February 15, 2024, \url{https://pedrocarvalho.es/docs/How%20Does%20WSPR%20Detect%20Aircraft%20over%20Long%20Distances%20-%2015FEB2024.pdf}; pp.~1--2, 5, 10 and 14--16.
\bibitem{oefgrpdf} R.~Godfrey and H.~Coetzee, Flight OE-FGR GDTAAA WSPRnet analysis: case study, October 28, 2022 (updated October 30, 2022), \url{https://www.dropbox.com/s/hlfi1amubynkjvo/Flight%20OE-FGR%20GDTAAA%20WSPRnet%20Analysis%20-%20Case%20Study%20-%2030th%20October%202022.pdf?dl=0}; pp.~9--10.
\bibitem{qtr901pdf} R.~Godfrey, H.~Coetzee and S.~Maskell, Flight QTR901 GDTAAA WSPRnet analysis: case study, June 8, 2023, \url{https://www.dropbox.com/s/224qont4xrgv54u/Boeing%20777-300ER%20QTR901%20PER%20DOH%2001NOV2022%20Detection%20and%20Tracking%20-%208th%20June%202023.pdf?dl=0}; pp.~10, 17, 27 and 32.
\bibitem{oefgr} R.~Godfrey, OE-FGR case study, \emph{The Search for MH370}, October 28, 2022, author's comments of May 18 and July 9, 2024, \url{https://www.mh370search.com/2022/10/28/oe-fgr-case-study/}.
\bibitem{wsprradar} R.~Godfrey, WSPR as radar, \emph{The Search for MH370}, December 18, 2023, \url{https://www.mh370search.com/2023/12/18/wspr-as-radar/}.
\bibitem{iannello2021} V.~Iannello, WSPR can't find MH370, \emph{MH370 and Other Investigations}, December 19, 2021, with comments by M.~Gilbert (December 21, 2021) and others, \url{https://mh370.radiantphysics.com/2021/12/19/wspr-cant-find-mh370/}.
\bibitem{kent} S.~Kent (signed sk999), WSPR (Weak Signal Propagation Reporting) and MH370: theory, April 11, 2021, updated September 14, 2021, \url{https://www.dropbox.com/scl/fi/oh0z7y6qipagrfr0b5142/Steve-Kent-WSPR-Paper.pdf?rlkey=3nl5lrt0yjos51ai8f1cl7x6k&dl=0}; p.~19.
\bibitem{loi} Loi Yi Yang Caden, Neo Jun Hui Gavin, Zhang Rui Xue and Ng Yuting, \emph{Forward Scatter Aircraft Detection with Amateur Radio Network WSPRnet}, January 21, 2025, \url{https://www.dsta.gov.sg/staticfile/ydsp/projects/files/reports/Report%20-%20Forward%20Scatter%20Aircraft%20Detection%20with%20Amateur%20Radio%20Network%20WSPRnet.pdf}, student research report, Raffles Institution and DSO National Laboratories, hosted by DSTA, Singapore, Sections I, III.B, IV.A and V.
\bibitem{liverpool} University of Liverpool, Researchers provide statistical expertise to help locate missing flight MH370, March 6, 2024, \url{https://news.liverpool.ac.uk/2024/03/06/university-researchers-provide-statistical-expertise-to-help-locate-mh370/}.
\bibitem{liverpooldesign} R.~Godfrey, comment of March 17, 2024, WSPR aircraft tracking, \url{https://www.mh370search.com/2024/02/15/wspr-aircraft-tracking/comment-page-1/#comment-2553}.
\bibitem{martencontrol} B.~Marten, WSPR Tracker, flight and control windows, \url{https://github.com/benmarten/wspr-tracker/blob/edd82f5d7c400dbae7f9b3f7e282bc2494859ecc/app/components/ControlComparison.vue}.
\bibitem{marten} Ben Marten, \emph{WSPR Tracker}, backtest of February 17, 2026, \url{https://github.com/benmarten/wspr-tracker/blob/edd82f5d7c400dbae7f9b3f7e282bc2494859ecc/data/tuning-recommendation.tiny.json}, \texttt{testMetrics.meanMedianKm}.
\bibitem{taylor} J.~Taylor and B.~Walker, WSPRing around the world, \emph{QST}, November 2010, 30--32, \url{https://wsjt.sourceforge.io/WSPR_QST_Nov_2010.pdf}.
\bibitem{westphal2021} R.~Westphal, ``Geocaching'' in the ionosphere, HamSCI Workshop 2021, March 2021, \url{https://hamsci.org/publications/2021/geocaching-in-the-ionosphere}; Introduction.
\bibitem{wsprlive} wspr.live, a public mirror of the WSPRnet database, \url{https://wspr.live/}.
\end{thebibliography}

\begin{appendices}
\section{Exact definitions}\label{app:defs}

\paragraph{Records and keys.} Exact duplicate records are dropped first. A record is then keyed by its reported time (to the second), band, transmitter and receiver; a key with several differing records is dropped entirely. One remaining record is one link-slot. Its slot is the two-minute interval of its period that contains the reported time, so a record timed on an odd minute falls in the slot that contains it. A link-slot is scored when it lies in a test slot, both locators have six characters, the path is at least 100~km long and its window holds at least five records.

\paragraph{Crossing labels.} The three snapshots of a slot are those at its nominal start and one and two minutes later. A snapshot position's age is the snapshot time minus the aircraft's last position time; positions older than 60~s are dropped, and a stale report seen in two snapshots counts once. Between two reports a position is interpolated on the sphere; outside them it is projected along the reported ground speed and track for up to 60~s, or held for up to 15~s without them. Tracks reported only on the ground are excluded; a track counts as high (at least 20{,}000~ft barometric, FL200) for the whole slot if any of its three snapshots reports that, and a track with unknown altitude counts as airborne but not high. Positions are evaluated at $\tau=1,6,\ldots,111$~s after the first selected record time of the slot (which can be odd-minute). A candidate screen keeps tracks whose middle position (else any defined one) lies within 90~km of the arc's corridor. At each $\tau$ the distance is the perpendicular distance to the minor arc when the foot point lies between the stations, else the distance to the nearer station; a change of sign of the signed perpendicular distance (which side of the arc) between successive $\tau$, at a point between the stations and with an along-arc step under 0.1~rad, sets the distance to zero. Distances are capped at 30~km. Crossed: high-track distance $\le5$~km. Eligible uncrossed: high-track distance $\ge25$~km and any-airborne distance $\ge5$~km. Both require labels (all three snapshots present) and a covered path.

\paragraph{Coverage.} For each period, a $1^\circ\times1^\circ$ cell is covered when some aircraft report, airborne or on the ground, falls in it in at least a quarter of the period's available snapshots. A path is covered when every cell containing one of its arc points, sampled at most 10~km apart with the ends included, is covered. Each label set uses the coverage of its own snapshots.

\paragraph{Designs.} Same-link matching: controls are eligible uncrossed link-slots of the same transmitter, receiver and band on the same date with $|t_{\text{control}}-t_{\text{case}}|\le1{,}800$~s in reported time. Shared-station matching: controls are eligible uncrossed link-slots with the same reported time and band that share the case's transmitter (as transmitter) or its receiver (as receiver), are not the same link, and have path length in $[0.8,1.25]$ times the case's. For each case, $d$ is its flag minus the mean flag of its controls; each $E$ is the sum of $d$ over its cases divided by its number of cases. All three label sets are matched separately.

\paragraph{Simulated increase.} For each date, independently with probability 0.10 and a seed fixed per date, each link-slot that is crossed under the real labels (inside the mask or not) gets 6~dB added to its SNR in a copy of the records. The SNR scores of every record are recomputed from the copy; frequency and drift scores are unchanged. The copy is analyzed with the real and the shifted labels.

\paragraph{Resampling and deletions.} The day bootstrap draws the dates with replacement, 10{,}000 times, and recomputes each $E$ from the drawn dates' sums. The block runs draw $\lceil n/L\rceil$ starting dates and take $L$ successive dates in date order from each, wrapping at the end, keeping the first $n$. The prespecified deletion drops one date's sums; the fuller one also replaces, for the two dates whose shifted positions came from the deleted date, that shifted label's sums by the date's other shifted label's.

\section{The anomaly rule as implemented}\label{app:rule}
Each line pairs a phrase of Godfrey's statement~\cite{oefgr} with what the code does. In the right column $t$ is the record's time, $n$ the number of records in its window, $x$ its SNR (dB) or frequency (Hz), $\bar x$ the window mean and $s$ the window's sample standard deviation.

\begin{center}
\small
\begin{tabular}{p{0.44\textwidth}p{0.48\textwidth}}
\toprule
Published rule & Implementation \\
\midrule
``a specific transmitterID-receiverID-frequency triple for a $\pm3$ hour analysis window'' & transmitter call sign, receiver call sign and band; the triple's records with time in $[t-3\,\mathrm{h},t+3\,\mathrm{h}]$, the record itself included \\
mean and standard deviation & $\bar x$ and $s$ (divisor $n-1$) \\
SNR or frequency ``not more than 0.75 standard deviation from the mean'' is discarded & $|x-\bar x|>0.75\,s$ flags; $s=0$ gives no flag; frequency in Hz \\
``green for non-zero drift''; ``discard drift anomalies where there is a continual drift'' & nonzero drift, and fewer than half the window's records with nonzero drift \\
``discard candidate SNR or Frequency anomalies based on a small count''; ``If the count is below five we extend up to $\pm3$ days'' & a record with $n<5$ gets no flag and is not scored; the $\pm3$-day extension is not implemented \\
mobile stations ``detected from their call signs or the fact that the Maidenhead Grid locator changes'' & records with a slash in a call sign, a transmitter call starting 0, 1 or Q, or a call sign with more than one locator string among all 28 downloaded hours (four- and six-character strings counted as different) are dropped \\
antenna positions from a separate database & centers of six-character locators (error up to about 5~km); four-character locators dropped \\
\bottomrule
\end{tabular}
\end{center}

A link-slot is flagged when any of its SNR, frequency or drift flags is set.

\section{Variants and sensitivity runs}\label{app:variants}
Each entry is $\Delta$ in percentage points, with its 95\% nominal day-bootstrap interval, except the mean-$|z|$ row: it is 100 times the same placebo-differenced matched contrast computed for the absolute standardized SNR score $|z|$ instead of the flag, so the two-point threshold does not apply to it. SD denotes standard deviation. FL200+ means a barometric altitude of at least 20{,}000~ft. The distance rows are half-open bins $[a,b)$ of the high-track distance, with the primary controls. The upper block uses the original analysis order for date resampling. Every lower-block row, including the unmodified baseline, and both kinds of deletion use dates in chronological order before resampling. A finite number of bootstrap draws can therefore give slightly different pooled limits for identical data. Only the primary and simulated-increase rows determine the outcome. The additional robustness condition is given in Section~\ref{sec:results}.

\begin{center}
\small
\begin{tabular}{lll}
\toprule
Pooled, 72 days & Same-link & Shared-station \\
\midrule
Primary (0.75 SD, FL200+, 5 km) & $+0.21$ ($-0.53$, $0.95$) & $+0.35$ ($-0.46$, $1.17$) \\
1.0 SD tier & $+0.33$ ($-0.60$, $1.19$) & $+0.22$ ($-0.69$, $1.10$) \\
SNR flag only & $+0.28$ ($-0.56$, $1.14$) & $+0.71$ ($-0.12$, $1.55$) \\
Frequency flag only & $-0.17$ ($-1.00$, $0.61$) & $-0.29$ ($-1.09$, $0.53$) \\
Drift flag only & $+0.25$ ($-0.00$, $0.49$) & $-0.14$ ($-0.53$, $0.23$) \\
Radius 1 nautical mile & $+0.26$ ($-0.48$, $0.98$) & $+0.46$ ($-0.48$, $1.38$) \\
Any airborne aircraft & $+0.25$ ($-0.45$, $0.92$) & $+0.16$ ($-0.64$, $0.95$) \\
No coverage mask & $-0.10$ ($-0.49$, $0.28$) & $-0.08$ ($-0.36$, $0.18$) \\
Distance 0--2 km & $+0.29$ ($-0.43$, $1.01$) & $+0.43$ ($-0.48$, $1.33$) \\
Distance 2--5 km & $-0.13$ ($-1.56$, $1.25$) & $+0.02$ ($-1.49$, $1.56$) \\
Distance 5--10 km & $+0.15$ ($-0.96$, $1.28$) & $-0.33$ ($-1.66$, $1.02$) \\
Distance 10--25 km & $+0.10$ ($-0.78$, $0.93$) & $+0.09$ ($-1.02$, $1.16$) \\
Mean $|z|$ of SNR, $\times100$ & $-0.11$ ($-1.10$, $0.82$) & $+0.04$ ($-1.13$, $1.24$) \\
Simulated $+6$~dB SNR, $p=0.10$ & $+2.96$ ($2.19$, $3.70$) & $+2.68$ ($1.77$, $3.60$) \\
\midrule
Unmodified, chronological order & $+0.21$ ($-0.52$, $0.94$) & $+0.35$ ($-0.43$, $1.15$) \\
Two-minute slot keys for matching & $+0.21$ ($-0.52$, $0.94$) & $+0.35$ ($-0.43$, $1.15$) \\
Odd-minute records dropped at matching & $+0.22$ ($-0.51$, $0.96$) & $+0.35$ ($-0.43$, $1.15$) \\
Strict control cutoffs & $+0.21$ ($-0.52$, $0.94$) & $+0.35$ ($-0.43$, $1.15$) \\
All three together & $+0.22$ ($-0.51$, $0.96$) & $+0.35$ ($-0.43$, $1.15$) \\
Blocks of 2 dates & $+0.21$ ($-0.56$, $0.99$) & $+0.35$ ($-0.48$, $1.15$) \\
Blocks of 3 dates & $+0.21$ ($-0.58$, $0.99$) & $+0.35$ ($-0.45$, $1.13$) \\
Blocks of 4 dates & $+0.21$ ($-0.55$, $0.96$) & $+0.35$ ($-0.39$, $1.05$) \\
Blocks of 6 dates & $+0.21$ ($-0.51$, $0.92$) & $+0.35$ ($-0.38$, $1.02$) \\
\bottomrule
\end{tabular}
\end{center}

\begin{center}
\small
\begin{tabular}{lll}
\toprule
Second sample, 36 days & Same-link & Shared-station \\
\midrule
Primary (0.75 SD, FL200+, 5 km) & $-0.29$ ($-1.22$, $0.65$) & $+0.46$ ($-0.64$, $1.65$) \\
1.0 SD tier & $-0.37$ ($-1.50$, $0.70$) & $+0.10$ ($-1.18$, $1.44$) \\
SNR flag only & $-0.58$ ($-1.48$, $0.28$) & $+1.21$ ($0.05$, $2.36$) \\
Frequency flag only & $-0.04$ ($-1.27$, $1.05$) & $-0.23$ ($-1.34$, $1.00$) \\
Drift flag only & $+0.07$ ($-0.26$, $0.38$) & $-0.06$ ($-0.66$, $0.47$) \\
Radius 1 nautical mile & $-0.30$ ($-1.21$, $0.57$) & $+0.41$ ($-0.85$, $1.72$) \\
Any airborne aircraft & $+0.20$ ($-0.67$, $1.02$) & $+0.14$ ($-0.95$, $1.28$) \\
No coverage mask & $-0.23$ ($-0.68$, $0.20$) & $+0.02$ ($-0.38$, $0.38$) \\
Distance 0--2 km & $-0.17$ ($-1.08$, $0.69$) & $+0.36$ ($-0.87$, $1.63$) \\
Distance 2--5 km & $-0.79$ ($-2.69$, $1.12$) & $+0.92$ ($-1.24$, $3.18$) \\
Distance 5--10 km & $-0.30$ ($-1.72$, $1.14$) & $+0.49$ ($-1.21$, $2.28$) \\
Distance 10--25 km & $+0.11$ ($-1.22$, $1.42$) & $+0.46$ ($-1.01$, $1.98$) \\
Mean $|z|$ of SNR, $\times100$ & $-0.93$ ($-2.13$, $0.19$) & $+0.47$ ($-1.18$, $2.05$) \\
Simulated $+6$~dB SNR, $p=0.10$ & $+2.29$ ($1.35$, $3.22$) & $+2.96$ ($1.75$, $4.25$) \\
\midrule
Unmodified, chronological order & $-0.29$ ($-1.22$, $0.65$) & $+0.46$ ($-0.64$, $1.65$) \\
Two-minute slot keys for matching & $-0.29$ ($-1.22$, $0.65$) & $+0.46$ ($-0.64$, $1.65$) \\
Odd-minute records dropped at matching & $-0.29$ ($-1.21$, $0.66$) & $+0.45$ ($-0.65$, $1.64$) \\
Strict control cutoffs & $-0.29$ ($-1.22$, $0.65$) & $+0.46$ ($-0.64$, $1.65$) \\
All three together & $-0.29$ ($-1.21$, $0.66$) & $+0.45$ ($-0.65$, $1.64$) \\
Blocks of 2 dates & $-0.29$ ($-1.21$, $0.62$) & $+0.46$ ($-0.68$, $1.57$) \\
Blocks of 3 dates & $-0.29$ ($-1.23$, $0.59$) & $+0.46$ ($-0.66$, $1.51$) \\
Blocks of 4 dates & $-0.29$ ($-1.23$, $0.50$) & $+0.46$ ($-0.53$, $1.45$) \\
Blocks of 6 dates & $-0.29$ ($-1.23$, $0.42$) & $+0.46$ ($-0.43$, $1.37$) \\
\bottomrule
\end{tabular}
\end{center}

\begin{center}
\small
\begin{tabular}{lll}
\toprule
First sample, 36 days & Same-link & Shared-station \\
\midrule
Primary (0.75 SD, FL200+, 5 km) & $+0.75$ ($-0.37$, $1.85$) & $+0.22$ ($-0.95$, $1.29$) \\
1.0 SD tier & $+1.08$ ($-0.26$, $2.33$) & $+0.34$ ($-0.86$, $1.37$) \\
SNR flag only & $+1.21$ ($-0.19$, $2.53$) & $+0.17$ ($-1.00$, $1.37$) \\
Frequency flag only & $-0.32$ ($-1.40$, $0.75$) & $-0.40$ ($-1.59$, $0.77$) \\
Drift flag only & $+0.43$ ($0.07$, $0.80$) & $-0.22$ ($-0.73$, $0.28$) \\
Radius 1 nautical mile & $+0.86$ ($-0.28$, $2.01$) & $+0.49$ ($-0.89$, $1.80$) \\
Any airborne aircraft & $+0.31$ ($-0.79$, $1.36$) & $+0.18$ ($-1.00$, $1.25$) \\
No coverage mask & $+0.04$ ($-0.62$, $0.72$) & $-0.20$ ($-0.57$, $0.17$) \\
Distance 0--2 km & $+0.79$ ($-0.36$, $1.94$) & $+0.48$ ($-0.86$, $1.74$) \\
Distance 2--5 km & $+0.57$ ($-1.51$, $2.60$) & $-0.98$ ($-3.07$, $0.97$) \\
Distance 5--10 km & $+0.63$ ($-1.12$, $2.37$) & $-1.23$ ($-3.27$, $0.67$) \\
Distance 10--25 km & $+0.08$ ($-1.01$, $1.15$) & $-0.31$ ($-2.01$, $1.26$) \\
Mean $|z|$ of SNR, $\times100$ & $+0.77$ ($-0.71$, $2.20$) & $-0.42$ ($-2.07$, $1.31$) \\
Simulated $+6$~dB SNR, $p=0.10$ & $+3.68$ ($2.52$, $4.82$) & $+2.34$ ($0.93$, $3.62$) \\
\midrule
Unmodified, chronological order & $+0.75$ ($-0.37$, $1.85$) & $+0.22$ ($-0.95$, $1.29$) \\
Two-minute slot keys for matching & $+0.75$ ($-0.37$, $1.85$) & $+0.21$ ($-0.95$, $1.29$) \\
Odd-minute records dropped at matching & $+0.77$ ($-0.34$, $1.87$) & $+0.22$ ($-0.94$, $1.30$) \\
Strict control cutoffs & $+0.75$ ($-0.37$, $1.85$) & $+0.22$ ($-0.95$, $1.29$) \\
All three together & $+0.77$ ($-0.34$, $1.87$) & $+0.22$ ($-0.94$, $1.30$) \\
Blocks of 2 dates & $+0.75$ ($-0.45$, $1.95$) & $+0.22$ ($-1.07$, $1.35$) \\
Blocks of 3 dates & $+0.75$ ($-0.49$, $1.97$) & $+0.22$ ($-1.04$, $1.34$) \\
Blocks of 4 dates & $+0.75$ ($-0.40$, $1.92$) & $+0.22$ ($-0.95$, $1.31$) \\
Blocks of 6 dates & $+0.75$ ($-0.32$, $1.81$) & $+0.22$ ($-1.01$, $1.22$) \\
Without December 1, 2025 (date only) & $+0.91$ ($-0.18$, $2.01$) & $+0.24$ ($-0.96$, $1.35$) \\
\bottomrule
\end{tabular}
\end{center}

The strict-cutoff runs equal their own unmodified chronological baselines in every sample. Dropping odd-minute reports changes only matching and leaves previously computed flags, distances and coverage intact. Without odd-minute records same-link matching has 36{,}237 (first sample), 38{,}323 (second) and 74{,}560 (pooled) cases and shared-station matching has 32{,}347, 34{,}868 and 67{,}215; with slot keys shared-station matching has 32{,}372, 34{,}882 and 67{,}254. Once odd-minute records are dropped, slot keys change no remaining time, so the three together equal the odd-minute run. Every lower-block run keeps both simulated-increase lower limits above zero. Over the prespecified single-date deletions the smallest is 0.72 (first sample), 1.16 (second) and 1.62 (pooled), and over the fuller deletions 0.56, 0.91 and 1.60.

In the first sample, seven of 36 fuller deletions put the same-link upper limit between 2.003 and 2.162 points. In the second sample, all 36 fuller deletions passed. Their largest upper limits were 0.88 points (same-link) and 1.999 points (shared-station, without February 1, 2021). The corresponding maxima for the prespecified second-sample deletions were 0.75 and 1.82 points.

\section{Sample dates}\label{app:dates}
\emph{First sample} (36 dates): December 1, 2022, and the first of each month from January 2023 to December 2025 except December 2023. A one-day pilot on December 1, 2023 (with November 1, 2023 as its shifted date), whose outcome statistic was not computed, set the design before the main analysis. In the pilot, with ``uncrossed'' meaning no aircraft within 50~km, only 41 of 175{,}839 covered link-slots were uncrossed.

\emph{Second sample} (36 dates): the 36 most recent first-of-month dates not used before, counting back from November 1, 2026, with a date missing any of its 488 ADS-B snapshots replaced by the next earlier unused complete date. December 1, 2020 (1 snapshot missing) and January 1, 2021 (4 missing) were replaced by October 1, 2020 and August 1, 2020; September 1, 2020 also lacked two. The sample is August, October and November 2020, February 2021 to November 2022, and January to November 2026.

In the second sample 96.0\% of covered link-slots are crossed and 48{,}557 are eligible uncrossed, against 96.8\% and 48{,}032 in the first. In each sample a date's shifted positions come from the next and the previous date of that sample, in date order, the last date wrapping to the first.
\section{Checks}\label{app:checks}
An independent spot checker, which recomputes flags and distances for sampled link-slots from the raw files, found one error after the first outcome had been seen: floating-point running sums lost the variance of small windows on 4 of the 36 days. The fix (exact integer window sums) changed 439{,}157 frequency scores and no SNR score, and moved $\Delta$ by 0.02 points (same-link $+0.73$ to $+0.75$, shared-station $+0.20$ to $+0.22$).

\paragraph{Spot checker.} The spot checker imports nothing from the pipeline. It recomputes, for sampled link-slots (half crossed, half eligible uncrossed by their real labels), the window size and the SNR, frequency and drift scores from the raw WSPR files, and the two aircraft distances from re-downloaded snapshots with different formulas and 1-second sampling, accepting distances within 1.5~km of the stored ones, or both at least 29.5~km (the 30~km cap). It does not check coverage, shifted labels, simulated scores or the estimators; it detects a distance changed to 30~km and an SNR score shifted by 3. After the fix it found no mismatch in 140 link-slots on 7 first-sample dates and in 60 on 3 second-sample dates.

\paragraph{Crossing geometry.} On 6 randomly drawn first-sample dates, a reference calculation that considers every aircraft in the raw snapshots and samples positions every second was compared with the main one for 960 link-slots near the cutoffs or used as controls. Three crossing or control labels change when the reference replaces the main calculation: two controls next to aircraft whose reported positions jump hundreds of kilometers in about a minute, and one odd-minute record evaluated at different sampling times. On the other 957 the distances agree within 0.5~km. This tests agreement on sampled records only.

\paragraph{Tallies.} The coverage, crossing and simulated-increase counts and the path-length medians of Section~\ref{sec:method} and the introduction come from \texttt{checks/tallies.py}, run on the units files; its output is \texttt{data/review/tallies.json}.

\paragraph{Exact flags.} \texttt{exactflags.py} recomputes, for every record of each date, the window sums in unbounded integers and flags $|z|>0.75$ by the exact comparison $16(nx-S_1)^2(n-1)>9n(nS_2-S_1^2)$, with no flag when $nS_2-S_1^2=0$, for the real and the simulated SNR and for the frequency offset. Over 45{,}033{,}331 scored link-slots no flag differs from the recorded run. On four dates (June, August, September and October 2025) a malformed report puts a frequency about $2^{32}$~Hz from its link's first value, beyond what 64-bit squares hold; the exact check covers these dates too.

\end{appendices}
\end{document}
